\section{同时驱动微环的两个方向}\label{sec:bidir}

第~\ref{sec:sweetspot}~节表明，单向微环里信号和失谐抖动作用在同一个量上：有效失谐。微环还有另一组共享这个失谐的模式，即反向传播的模式。
用同一台激光泵浦时，两个方向看到相同的失谐抖动和光源强度噪声，而信号可以放在两者的差上。本节检验这样做有没有用。

\subsection{模型}

每个方向各有自己的腔内场。反向传播的两个场每个往返相互穿过一次，所以彼此只通过对方的往返平均功率起作用（交叉相位调制，系数 2）：
\begin{equation}
\frac{\mathrm{d}E_X}{\mathrm{d}n} = \hat{L}E_X + i\gamma L\left(|E_X|^2 + 2\langle|E_Y|^2\rangle\right)E_X + \sqrt{\theta_1}\,E_{\mathrm{in},X},
\qquad (X,Y)\in\{(\mathrm{cw},\mathrm{ccw}),(\mathrm{ccw},\mathrm{cw})\},
\label{eq:bidir}
\end{equation}
其中 $E_X(\tau)$ 是方向 $X$ 的快时间场，$\hat L$ 对其傅里叶分量的作用与式~\eqref{eq:lle} 中的线性项（损耗、失谐、色散）相同，$\langle|E_Y|^2\rangle=t_R^{-1}\int_0^{t_R}|E_Y|^2\,\mathrm{d}\tau$ 是另一方向的往返平均功率；这一耦合导致方向性双稳态~\cite{kaplan1982}、Kerr 微腔中反向传播光的自发对称破缺~\cite{delbino2017}以及反向传播孤子~\cite{yang2017}。我们忽略背向散射。
失谐抖动和光源强度噪声对两个方向相同，量子场噪声和探测噪声相互独立。两个方向各从自己的下载端口输出（图~\ref{fig:bidir}a），共探测 $2\times31$ 根梳齿（$62\times8$ 个读出特征）。
用两个方向的方案把工作点的功率平分给两个方向。关掉反向场时，这个仿真与单向仿真器的结果完全一致（无噪声时输出相同）。

我们在加大下载端口耦合组的同一批随机工作点上比较：
\emph{单向}（读 31 或 62 根梳齿；后者与双向方案的特征数相同）；
\emph{双向、信号反相}（反向用同一个掩模，信号 $u\to1-u$）；
\emph{双向、只有一侧带信号}（反向用同一个掩模、符号恒定）；
\emph{双向、独立掩模}；以及耦合的对照\emph{两个独立微环、信号反相}（把反相方案的两路驱动分别送进两个不耦合、由同一台激光供光的微环：特征相同、共模噪声相同，只是没有交叉相位调制）。

\subsection{随机搜索}

表~\ref{tab:bidir}(a) 每种方案用 768 个随机工作点（各方案用同一批点），流程与表~\ref{tab:encodings} 相同：在第一对运行最好的 8 个点中用 3 对新运行选点，
再在 3 条序列上各用 3 对新运行报告，分有、无失谐抖动两种。起作用的是第二路的信号反相：这样驱动的两个独立微环，比读出同样多梳齿的单向微环好 4 倍；
同一个微环的两个方向再好 1.7 倍。只在一侧加信号没有好处（192 个点的小搜索中为 $2.0\times10^{-4}$，同一搜索的单向为 $1.6\times10^{-4}$），
两个方向用独立掩模也几乎没有好处（$1.4\times10^{-4}$）。好的双向工作点很少：768 个点中单对运行验证误差低于 $10^{-4}$ 的有 8 个（两个独立微环 4 个，单向 0 个）；
另外两次各 192 个点的小搜索中，双向方案分别找到 $2.4\times10^{-5}$ 和 $1.8\times10^{-4}$，两个独立微环则为 $8.0\times10^{-5}$ 和 $7.7\times10^{-5}$。

\begin{table}[htbp]
\centering
\caption{双向驱动，Mackey--Glass，现实噪声下有、无失谐抖动时的测试 NMSE（3 条序列的几何平均）。
(a) 每种方案 768 个随机工作点（同一批点）中最好的，流程同表~\ref{tab:encodings}。
(b) 固定工作点上的最优掩模：先随机搜索 512 个掩模，再做进化（6 代，每代在最好的 16 个周围产生 240 个），在序列 0--2 上选定，
在搜索和选择都没有用过的序列 3--5 上报告；“随机掩模”是该工作点原来的随机掩模在同样的留出序列上的结果。}
\label{tab:bidir}
\small\setlength{\tabcolsep}{4pt}
\begin{tabular}{lcccc}
\toprule
 & \multicolumn{2}{c}{(a) 随机搜索} & \multicolumn{2}{c}{(b) 优化掩模，留出序列} \\
\cmidrule(lr){2-3}\cmidrule(lr){4-5}
方案 & 有抖动 & 无抖动 & 随机掩模 & 优化后（无抖动） \\
\midrule
单向，31 根 & $1.7\times10^{-4}$ & $8.6\times10^{-5}$ & -- & -- \\
单向，62 根 & $1.4\times10^{-4}$ & $4.7\times10^{-5}$ & $1.9\times10^{-4}$ & $3.3\times10^{-5}$（$1.4\times10^{-5}$） \\
两个独立微环，信号反相 & $3.6\times10^{-5}$ & $1.7\times10^{-5}$ & $4.0\times10^{-5}$ & $1.1\times10^{-5}$（$5.3\times10^{-6}$） \\
双向，信号反相 & $\mathbf{2.1\times10^{-5}}$ & $\mathbf{9.4\times10^{-6}}$ & $1.7\times10^{-5}$ & $\mathbf{7.3\times10^{-6}}$（$4.4\times10^{-6}$） \\
\bottomrule
\end{tabular}
\end{table}

好处并不来自抵消共模抖动：每种方案的抖动损失（有抖动与无抖动误差之比）都在 2 到 3 之间。好处来自计算本身更好，即无抖动时的误差更低。
一个可能的机制是式~\eqref{eq:bidir} 中的交叉相位项：它让一个方向每根梳齿的失谐随另一方向的往返平均功率移动，而信号反相时，这个功率对输入的响应符号相反。
于是信号经两条符号相反的路径作用在有效失谐上，形成一个全局的、随输入变化的门控项，并带来两个方向特征之间的乘积。我们没有进一步拆分这一点。

\subsection{稳健性与掩模的作用}

在第一次搜索得到的最好双向工作点周围（表~\ref{tab:bidir}(a) 用的是大搜索中的另一个点），我们每次只改一个参数，每个变体在 3 条序列上各用 2 对新运行
（该点本身为 $2.5\times10^{-5}$）。归一化失谐在 1.5 到 2.1 之间（超过 2.2 后急剧升到约 $10^{-3}$）、功率为原来的 0.8 到 1.1 倍、调制深度不低于原来的 0.9 倍、
边带功率比例变化 $\pm0.2$、输入耦合 $\log_{10}\Qone$ 在原值 $-0.1$ 到 $+0.2$ 之间时，误差都在 $1.7\times10^{-5}$ 到 $3.8\times10^{-5}$ 之间。
所以这个工作点是一片宽的区域，不是一个窄尖峰。但掩模不能随便换：物理参数不变，换 16 个别的随机掩模，误差在 $9.8\times10^{-5}$ 到 $8.7\times10^{-3}$ 之间。
这表明好的工作点稀少，主要是因为好的掩模稀少。

掩模由可编程谱整形器写入，可以像训练参数一样选择。表~\ref{tab:bidir}(b) 给每种方案在其最好工作点上同样的优化预算（双向方案的失谐移到上述平台内的 1.73）。
优化对每种方案都有效，在留出序列上好 2.3 到 5.7 倍，并缩小了方案之间的差距：两个独立微环比单向好 2.9 倍，同一微环的两个方向再好 1.55 倍，刚好在我们视为持平的门槛上。
最好的结果为有抖动时 $7.3\times10^{-6}$，比第~\ref{sec:readout}~节最好的单环（$1.55\times10^{-4}$）低 21 倍。这个倍数包含了好几项改变，不应理解为单单第二个方向的作用：读出加倍（62 根而不是 31 根）、掩模经过优化、驱动了第二个方向，而且是在留出序列上用较少的运行对评估的。同条件比较时，两个方向比读出规模相同、掩模同样优化的单向微环好 4.5 倍，比两个独立微环好 1.55 倍。

\subsection{输入延迟线：把时间搬到频率上}

让各输入梳齿带上不同延迟的信号，就能把过去的几个符号同时送进微环，四波混频由此形成跨时刻的乘积（第~\ref{sec:kernel}~节）。
我们把 $|\mu|=3$--4、5--6、7--8 的梳齿分别延迟 1、2、3 个步长（每步 $s$ 个符号；$|\mu|=1$--2 不延迟），重做随机搜索（192 个点，各方案同一批点）。
单向（62 根）的误差从 $2.4\times10^{-4}$（$s=0$）降到 $1.0\times10^{-4}$（$s=1$）和 $6.1\times10^{-5}$（$s=2$）；
双向信号反相只从 $2.1\times10^{-5}$ 降到 $1.6\times10^{-5}$ 和 $1.7\times10^{-5}$。延迟线补上了单向微环缺的很大一部分，对双向方案帮助不大。

\subsection{与同等优化的软件比较}\label{sec:bidir-software}

公平的比较要给软件同样的优化。我们用同样的流程和预算，在四个输入噪声水平下优化 62 节点 ESN 的输入权重和偏置（它们对应掩模；循环矩阵固定不动，就像微环的物理不变），
并在同样的留出序列上报告。优化后的 ESN 在输入噪声 $10^{-4}$、$3\times10^{-4}$、$10^{-3}$、$3\times10^{-3}$ 下分别为 $1.7\times10^{-7}$、$6.6\times10^{-7}$、
$6.5\times10^{-6}$、$4.5\times10^{-5}$（随机输入权重时为 $2.5\times10^{-6}$、$5.5\times10^{-6}$、$2.2\times10^{-5}$、$8.7\times10^{-5}$）。
因此最好的双向微环（$7.3\times10^{-6}$）相当于读出维度相同、经过同等优化、输入噪声约 $10^{-3}$ 的 ESN，按第~\ref{sec:software}~节的换算约为 8 位，
比随机掩模的最好单环多 1 位。在这一速率下转换器的 9 到 10 位精度处，ESN 好 3 到 11 倍。双向驱动、掩模优化和延迟线缩小了与同规模软件的差距，但没有消除它（图~\ref{fig:bidir}b）。


\begin{figure}[htbp]
\centering
\includegraphics[width=\linewidth]{fig10_bidirectional}
\caption{双向驱动。(a) 示意图：同一台激光和输入光梳送出两路掩模，分别携带信号 $u$ 和反相信号 $1-u$，进入同一微环的顺时针（cw）和逆时针（ccw）方向，
两个方向通过交叉相位调制（XPM）耦合；每个方向从自己的下载端口输出，各探测 31 根梳齿。对照组把同样的两路送进两个独立微环。
(b) 留出序列 3--5 上的测试误差随 62 节点 ESN 输入噪声的变化，虚线为随机输入权重，实线为优化后的输入权重；水平线为优化后的单向微环（62 根）、
两个独立微环和双向微环（表~\ref{tab:bidir}(b)，有抖动）。上方横轴为等效转换器位数，$2^{-N}/\sqrt{12}$。}
\label{fig:bidir}
\end{figure}
